% **************************************************************************
%           Testsysteme
% **************************************************************************
\subsection{Definition Legacy-Code}
\label{ch:SdF:Legacy:sec:Legacy_ungetestet}
% Gründe zum Ändern von Software > Testsystemen

Soll Software verändert werden, gibt es vier mögliche Gründe dafür. Entwickler können ein Feature hinzufügen, einen Fehler beheben, ein Refactoring oder eine Optimierung vornehmen. % oder das Verhalten in Form von Speicher oder Ressourcen optimieren.
Dies impliziert die Veränderung der Code-Struktur oder bestehender Funktionalität, der Einbau von neuer Funktionalität oder die Anpassung des Ressourcen-Verbrauchs. Dabei soll das bestehende Verhalten bewahrt werden. Das bedeutet, dass der Code so wenig wie möglich aber so viel wie nötig verändert werden soll \cite[S. 7]{Feathers:2013}. % S. 7

% Legacy > ohne Tests -> Verhalten bewahren = Testen -> Relevanz von Tests
Während eine aufgeräumte und gepflegte Umgebung dem Entwickler hilft, Überblick über den Code zu behalten, reicht dies in wachsenden und vor allem Legacy-Systemen nicht dafür aus, alle Einzelheiten zu kennen. Wird an dem Code eine Änderung vorgenommen, ist es kaum möglich, zu erkennen, ob der Code dadurch verbessert oder verschlechtert wurde. Zudem ist es ineffizient für den Entwickler, jede Einzelheit per Hand zu testen, ob sich das Verhalten exakt so verändert hat, wie es vorgesehen war. \glqq Code ohne Tests ist schlechter Code\grqq{} \cite[S. XVI]{Feathers:2013}. El-Morabea et al. ergänzt, dass ungetesteter Code wie eine tickende Zeitbombe sei und jederzeit in Form von Bugs oder Abstürzen explodieren könne \cite[S. 17]{ElMorabea:2021}.
Um zu garantieren, dass sich an dem Verhalten des bestehenden Codes bei Codeänderungen nichts verändert, werden Tests benötigt \cites[S. XVI]{Feathers:2013}[S. 128]{ElMorabea:2021}[S. 25]{Oppel:2025}. % S. XVI

Tests sind ein \glqq wesentlicher Bestandteil des Entwicklungszyklus\grqq{} \cite[S. 16]{Oppel:2025}. Sie helfen Entwicklern, hohe Qualitätsstandards einzuhalten und garantieren die Stabilität und Zuverlässigkeit von Softwareprodukten \cite[S. 25]{Oppel:2025}.
Wird die Software ausreichend getestet, sieht der Entwickler sofort, wenn sich das Code-Verhalten ungewollt verändert hat. Dafür wird der Aufbau und die Pflege eines umfangreichen Test-Frameworks zur regelmäßigen und häufigen Überprüfung von der einzelnen Funktionseinheit bis zum komplexen Zusammenspiel verschiedener Services benötigt \cites[S. 45]{Feathers:2013}[S. 3]{ElMorabea:2021}.

Zur Reduktion des manuellen Aufwands werden diese Testframeworks automatisiert. Dieses hilft weiterhin zur Steigerung der Testabdeckung und ermöglicht dem Entwickler eine direkte Rückmeldung. Da sich somit auch der Arbeitsprozess effizienter gestaltet, kann das Unternehmen von Kosteneinsparungen profitieren \cite[S. 24]{Oppel:2025}.

Mithilfe einer solchen Testbasis lässt sich Code schnell modifizieren, ohne dass der Entwickler den vollständigen Überblick benötigt.
Zudem werden auch strukturelle Veränderungen (\textit{Refactoring}), welche das Verhalten nicht verändern, durch die Testunterstützung vereinfacht \cites[S. XVI, 5, 45]{Feathers:2013}[S. 128]{ElMorabea:2021}.   % [S. XVI, 5, 45]
